\section{继续预训练与阶段检查}
\label{sec:nlp-llm-continuation}

训练预算用完，只说明过程到达了预定终点。是否继续学习、采用哪组参数以及下一阶段还缺什么，需要依据未参与更新的数据和明确的能力问题判断。这里讨论预训练阶段内部的检查与适应；怎样构造完整评价协议、分析统计不确定性和污染风险，将在\ref{chap:nlp-evaluation}章统一展开。

\subsection{训练进展与参数选择}
\label{sec:nlp-llm-checkpoints}

训练损失反映当前采样批次的拟合，留出损失反映模型在同一约定下对未用于更新材料的预测。二者应采用相同的词元化、可见范围和计分方式；若留出材料的来源配比不同，还需分别报告来源结果。否则，整体损失变化可能仅来自比较对象变化，而非模型对同一分布预测得更好。

参数检查点是在训练过程中保存的参数及必要状态。选择检查点时，可以观察总体留出损失，并检查关心的来源和能力。若总体损失下降但某种少数语言持续恶化，按一个总均值选择参数就会掩盖该语言的退步。若多次根据留出结果修改配比和超参数，这批留出数据已经参与选择，最终结论还需要未反复用于决策的测试材料。

曲线平台期有多种可能解释：模型接近当前目标下的可达状态，学习率过小，材料过度重复，或者目标难以充分提供所需能力的信号。单凭曲线变平不能区分这些原因。应结合数据曝光、更新范数、各来源损失及异常记录检查，再决定增加数据、改变日程还是停止。突然的损失尖峰也应先排查样本与更新，而不能直接称为能力突然下降。

能力检查需要与本阶段目标相符。预训练模型可能正确续写知识性文本，却在直接问句后继续写另一个问题，因为它尚未被要求稳定扮演助手。把这种行为问题与知识缺失区分后，才能判断是需要更多预训练材料，还是应该交给有监督微调。检查的作用是定位下一步应改变的环节，不是用一个总榜单代替训练诊断。

\subsection{新领域与新语言的继续学习}
\label{sec:nlp-llm-domain-continuation}

若模型已经具有一般语言规律，但专业术语或某种语言覆盖不足，可以从现有参数出发，改用相关材料继续优化预训练目标。这称为\term{继续预训练}（continued pre-training）：仍从文本本身构造预测信号，主要变化在训练分布、曝光比例和更新安排。它与下一章根据指令—回答示范学习行为不同，虽然两种过程都从已有参数开始。

设原有材料分布为$P_{\mathrm{old}}$，新材料分布为$P_{\mathrm{new}}$，可以用$P_\rho=(1-\rho)P_{\mathrm{old}}+\rho P_{\mathrm{new}}$控制两者混合，$\rho\in[0,1]$。增大$\rho$使同一预算内更多更新针对新材料，但不保证旧分布目标同时下降。为解释这个冲突，设两类损失在当前参数的梯度分别为$\bm g_{\mathrm{old}}$和$\bm g_{\mathrm{new}}$。只沿新材料更新$\Delta\theta=-\eta\bm g_{\mathrm{new}}$时，旧损失的一阶变化近似为
\begin{equation}
 \Delta\mathcal L_{\mathrm{old}}
 \approx-\eta\langle\bm g_{\mathrm{old}},\bm g_{\mathrm{new}}\rangle.
 \label{eq:nlp-llm-continuation-conflict}
\end{equation}
若内积为负，这一步可能增加旧损失。该式是局部近似，解释目标冲突为何可能发生，不提供整个训练过程的变化保证。

保留一定旧材料、缩小更新强度或限制继续学习预算，都是控制改变范围的方式，但其具体作用需要新旧数据共同检查。对材料很少的新语言，单纯反复复制同一语料可能提高训练拟合，却增加记忆风险。词元化也影响条件：若新文字被切成极多字节片段，相同最大长度能容纳的语义内容会减少；扩展词表则还需要学习新增嵌入与输出参数，不能只改编码表而沿用原接口假设。

继续学习还需要区分事实更新与能力退步。活动从周四改到周五，模型对旧安排降低概率可能是正确更新；但因此失去读取其他活动时间的能力才是另一类问题。比较时应固定问题所依据的时间和材料，明确期望保留的是旧事实、通用技能还是输入协议，避免把一切行为变化都称为遗忘。

\subsection{上下文长度的继续适应}
\label{sec:nlp-llm-context-extension}

网络在公式上能够接收更长序列，不意味着它已经学会使用更远的证据。若训练中距离最大为$L$，使用时出现$sL$的序列，模型同时面对新位置范围、新距离组合和更多干扰内容。这里需要把位置机制的改变与长序列学习一起考虑，而不能只提高一个最大长度配置。

以式\eqref{eq:nlp-llm-rope}为例，直接把位置$t$扩展到原来范围之外，会引入训练中没有出现过的旋转组合。位置插值（positional interpolation）用$t/s$替代$t$，把更长序列的位置压回原有范围。若原最大位置为$L$、目标扩展倍数为$s>1$，则新序列末端$sL$映射到$L$附近；相邻位置的间隔也从1缩为$1/s$，所以局部相对关系同时发生变化。这解释了为什么插值可以限制位置范围，却仍需要通过继续学习适应新的距离分布\scite{nlp-llm-pi2023}。

YaRN采用随频率变化的缩放，综合保留部分原频率与插值后的频率，并调整注意力分数尺度，以减少统一压缩对不同频段的影响\scite{nlp-llm-yarn2023}。它是位置处理与学习安排的组合，不能仅凭频率公式保证任意长度的正确检索。评价时需要区分训练过的长度、插值范围以及进一步外推的长度；超过继续学习范围的结果不能与范围内结果合并成一个“支持长度”。

长序列样本也必须真正包含跨距离的预测需要。把大量互不相关短文档串成一条很长记录，虽然增加输入长度，却未必要求模型使用远距离关系。可以检查在不同距离放置改期通知、随后询问最终安排时，模型是否使用了正确证据，以及加入相似但无关安排后怎样变化。局部预测损失、跨段检索和长材料综合各考察不同问题，单个针式检索例子不能代表所有长文本能力。

\subsection{原有能力保持与阶段交接}
\label{sec:nlp-llm-pretraining-handoff}

完成初始或继续预训练后，应把新参数放回相同协议下与旧参数比较。对目标领域，检查相同留出材料上的损失和所需任务结果；对原有范围，检查语言覆盖、基本知识、长短上下文和必要行为边界。比较既要看变化量，也要看变化后的水平。例如，一个错误率从80\%降到60\%的模型虽然下降二十个百分点，仍不能据此称为已满足低错误率要求。

检查还应覆盖词元与位置接口。模型参数、词元化规则、特殊标记、结束符号和上下文配置需要作为一组产物保存；把不匹配的聊天模板或词表接到参数上，会产生与模型学习能力不同的运行问题。数据来源、时间范围、混合比例和训练目标说明，则帮助后续阶段判断哪些能力已经接触过、哪些材料仍不足。

进入有监督微调并不要求预训练模型已经具有完整助手行为，但需要知道其起点和缺口。若主要问题是对新领域语言缺少基本概率建模，少量交互示范未必足够补齐；若内容知识已经存在而输出方式不稳定，示范学习才直接针对这个行为差异。\ref{chap:nlp-supervised-finetuning}章接住的是这个明确的起点，\ref{chap:nlp-reasoning}章则把推理所需材料、轨迹与反馈放回多个学习阶段综合讨论。
